\documentclass[11pt,letterpaper]{article}

% ============================================================
% TIPOGRAFÍA — COMPILAR CON XeLaTeX
% ============================================================
\usepackage{fontspec}
\setmainfont{Andika}

% ============================================================
% MÁRGENES
% ============================================================
\usepackage[margin=1.6cm]{geometry}

% ============================================================
% MATEMÁTICAS
% ============================================================
\usepackage{amsmath,amssymb,amsthm}

% ============================================================
% TABLAS, GRÁFICOS Y COLOR
% ============================================================
\usepackage{graphicx}
\usepackage{float}
\usepackage{booktabs}
\usepackage{xcolor}
\usepackage{array}

% ============================================================
% DOS COLUMNAS
% ============================================================
\usepackage{multicol}
\setlength{\columnsep}{18pt}
\setlength{\columnseprule}{0.2pt}

% ============================================================
% CONFIGURACIÓN GENERAL
% ============================================================
\setlength{\parindent}{0pt}
\setlength{\parskip}{3pt}
\linespread{1.02}

% Numeración simple: 1. 2. 3. ...
\renewcommand{\theenumi}{\arabic{enumi}}
\renewcommand{\labelenumi}{\theenumi.}

% Listas compactas sin enumitem
\makeatletter
\renewenvironment{enumerate}
  {\ifnum \@enumdepth >\thr@@\@toodeep\else
   \advance\@enumdepth\@ne
   \edef\@enumctr{enum\romannumeral\the\@enumdepth}%
   \expandafter\list\csname label\@enumctr\endcsname
     {\usecounter\@enumctr
      \setlength{\topsep}{3pt}%
      \setlength{\itemsep}{2pt}%
      \setlength{\parsep}{0pt}%
      \setlength{\partopsep}{0pt}%
      \setlength{\leftmargin}{1.6em}%
      \setlength{\labelsep}{0.4em}}%
   \fi}
  {\endlist}
\makeatother

\newcommand{\problema}[1]{%
  \vspace{2pt}
  \noindent\textbf{\large PROBLEMA #1}\par
  \vspace{2pt}
}

\begin{document}
\pagestyle{empty}

\begin{multicols}{2}
\raggedcolumns

% ============================================================
% PROBLEMA 1
% ============================================================
\problema{1}

Un camión de entregas viaja de una ciudad A a otra ciudad B. El camión puede
presentar entre 0 y 3 inconvenientes (averías o accidentes) en sus recorridos
de ida y regreso. Sea \(X\) el número de inconvenientes en el viaje de ida y
\(Y\) el número de inconvenientes en el viaje de regreso. La distribución
conjunta de \(X\) y \(Y\) está dada por la siguiente tabla:

\begin{center}
\small
\setlength{\tabcolsep}{5pt}
\begin{tabular}{c|cccc}
\toprule
\(f(x,y)\) & \(y=0\) & \(y=1\) & \(y=2\) & \(y=3\) \\
\midrule
\(x=0\) & 0.01 & 0.02 & 0.07 & 0.01 \\
\(x=1\) & 0.03 & 0.06 & 0.10 & 0.06 \\
\(x=2\) & 0.05 & 0.12 & 0.15 & 0.08 \\
\(x=3\) & 0.02 & 0.09 & 0.08 & 0.05 \\
\bottomrule
\end{tabular}
\end{center}

\begin{enumerate}
\item Cuál es la probabilidad de que un camión pueda hacer el recorrido completo sin inconvenientes?
\item Cuál es la probabilidad de que un camión tenga un inconveniente en el viaje de ida?
\item Cuál es la probabilidad de que un camión encuentre dos inconvenientes en el viaje de regreso?
\item Cuál es la probabilidad de que presente menos de cuatro inconvenientes en todo el recorrido (ida y regreso)?
\item En cien viajes con tres inconvenientes en el recorrido de regreso, cuántos se espera que presenten dos inconvenientes en el viaje de ida?
\item Determinar las funciones de probabilidad marginales de \(X\) y \(Y\).
\item Determinar la función de probabilidad condicional \(P(Y\mid X=0)\).
\item Cuántos inconvenientes se esperan en un viaje completo?
\item Construir los gráficos de la función de probabilidad conjunta, de las marginales y de la condicional.
\item Se pueden considerar independientes las variables \(X\) y \(Y\)?
\end{enumerate}

% ============================================================
% PROBLEMA 2
% ============================================================
\problema{2}

Sea \(X\) la cantidad de encogimiento (\%) de una fibra cuando se calienta a
\(120\,^{\circ}\mathrm{C}\), y \(Y\) el encogimiento adicional cuando se
calienta a \(140\,^{\circ}\mathrm{C}\). La función de densidad conjunta de
\(X\) y \(Y\) está dada por

\[
f(x,y)=
\begin{cases}
kxy^2, & 3\leq x\leq4,\quad 0.5\leq y\leq1,\\
0, & \text{en otro caso}.
\end{cases}
\]

\begin{enumerate}
\item Determinar el valor de \(k\) para que \(f(x,y)\) sea una función de densidad válida.
\item Qué porcentaje de las fibras tienen un encogimiento menor al \(3.2\%\) a \(120\,^{\circ}\mathrm{C}\) y un encogimiento adicional mayor al \(0.8\%\) a \(140\,^{\circ}\mathrm{C}\)?
\item Si 500 fibras presentaron un encogimiento adicional menor al \(0.8\%\) a \(140\,^{\circ}\mathrm{C}\), cuántas de ellas tuvieron un encogimiento menor al \(3.8\%\) a \(120\,^{\circ}\mathrm{C}\)?
\item Son independientes las variables \(X\) y \(Y\)?
\end{enumerate}

% ============================================================
% PROBLEMA 3
% ============================================================
\problema{3}

Una empresa produce mezclas de café con tres variedades: Colombia, Bourbon y
Caturra. Sea \(X\) la proporción del peso correspondiente a la variedad
Colombia y \(Y\) la proporción del peso correspondiente a la variedad Bourbon
en un paquete de café. La función de densidad conjunta de \(X\) y \(Y\) está
dada por

\[
f(x,y)=
\begin{cases}
24xy, & 0\leq x\leq1,\quad 0\leq y\leq1,\quad x+y\leq1,\\
0, & \text{en otro caso}.
\end{cases}
\]

\begin{enumerate}
\item Qué porcentaje de los paquetes contiene más de la mitad de su peso en café tipo Colombia?
\item Si se distribuyeron 200 paquetes con menos de la mitad de su peso en Bourbon, cuántos de ellos contienen más de \(3/4\) de su peso en Colombia?
\item Si 300 paquetes contienen la mitad de su peso en Colombia, cuántos de ellos contienen menos de la mitad de su peso en Bourbon?
\item Cuál es la probabilidad de que el café Bourbon represente más de la mitad de la mezcla?
\item Determinar las funciones marginales de \(X\) y \(Y\).
\item Graficar las funciones marginales \(g(x)\) y \(h(y)\).
\end{enumerate}

\end{multicols}
\end{document}
